\documentclass[11pt,a4paper]{article}
\usepackage[utf8]{inputenc}
\usepackage[margin=1in]{geometry}
\usepackage{amsmath,amssymb}
\usepackage{graphicx}
\usepackage{booktabs}
\usepackage{hyperref}
\usepackage{enumitem}
\usepackage{xcolor}
\usepackage{algorithm}
\usepackage{algpseudocode}

\title{Discrete Recursive Models for Continuous Robotic Control:\\
Applying TRM and HRM to PushT}
\author{VLA-HRM Project}
\date{\today}

\begin{document}
\maketitle

\begin{abstract}
We investigate applying discrete recursive reasoning models --- Tiny Recursive Models (TRM) and Hierarchical Reasoning Models (HRM) --- to the PushT robotic manipulation benchmark from Diffusion Policy. These models, originally designed for discrete puzzles (Sudoku, Mazes, ARC-AGI), perform iterative refinement through recursive computation. We bridge the continuous-discrete gap through careful tokenization of action and observation spaces, and introduce dimension-aware positional encoding to preserve the structure of flattened multi-dimensional sequences. We conduct extensive experiments on tokenization strategies, positional encodings, model sizes, and recursion depths.
\end{abstract}

\section{Introduction}

Diffusion Policy~\cite{chi2023diffusion} has shown strong performance on robotic manipulation tasks by formulating action prediction as a denoising diffusion process. However, recent work on discrete recursive models --- TRM~\cite{trm2024} and HRM~\cite{hrm2024} --- has demonstrated remarkable reasoning capabilities on tasks like Sudoku and ARC-AGI with very compact models (7--27M parameters).

A natural question arises: \textit{Can the recursive reasoning mechanism of these models be leveraged for robotic control?} The iterative refinement process in TRM/HRM is analogous to the denoising process in diffusion models --- both progressively transform an initial estimate into a final prediction. However, the key challenge is that robotic actions are continuous, while TRM/HRM operate on discrete tokens.

\subsection{Contributions}
\begin{itemize}[noitemsep]
    \item Adaptation of TRM and HRM architectures for continuous robotic control via discrete tokenization
    \item Dimension-aware positional encoding for structured multi-dimensional sequences
    \item Extensive ablation studies on tokenization, positional encoding, model size, and recursion depth
    \item Comparison against autoregressive GPT baseline and Diffusion Policy
\end{itemize}

\section{Background}

\subsection{PushT Task}
PushT is a planar pushing task where a circular agent must push a T-shaped block to match a target pose. The state space is 5-dimensional: agent position $(x_a, y_a)$, block position $(x_b, y_b)$, and block angle $\theta_b$. Actions are 2-dimensional target positions $(x_t, y_t)$ for PD control. The evaluation metric is coverage (intersection area / goal area), with success at 95\% coverage.

The dataset contains 206 human demonstrations with 25,650 total timesteps. Following Diffusion Policy conventions, we use an observation horizon of $T_o = 2$ and a prediction horizon of $T_a = 16$.

\subsection{Recursive Reasoning Models}

\paragraph{TRM (Tiny Recursive Model)} uses a single reasoning module that is recursively applied in a two-level loop:
\begin{itemize}[noitemsep]
    \item \textbf{L-level}: Low-level computation cycles ($L_{\text{cycles}}$ iterations)
    \item \textbf{H-level}: High-level refinement cycles ($H_{\text{cycles}}$ iterations)
\end{itemize}
The same weight-shared module processes both levels, alternating between updating the low-level state $z_L$ (with high-level input injection) and the high-level state $z_H$ (with low-level input injection).

\paragraph{HRM (Hierarchical Reasoning Model)} extends TRM with \textit{separate} modules for the H-level and L-level:
\begin{itemize}[noitemsep]
    \item H-level module: slow, abstract planning with its own transformer layers
    \item L-level module: fast, detailed computation with its own transformer layers
\end{itemize}

Both models only backpropagate through the final reasoning cycle, making training memory-efficient despite deep effective computation.

\section{Method}

\subsection{Tokenization}

We discretize continuous observations and actions into a vocabulary of $V$ bins. We explore two strategies:

\paragraph{Uniform Binning} divides each dimension's range $[\text{min}, \text{max}]$ into $V$ equal-width bins:
\begin{equation}
    \text{token}(x) = \left\lfloor \frac{x - x_{\min}}{x_{\max} - x_{\min}} \cdot V \right\rfloor, \quad
    \hat{x}(\text{token}) = x_{\min} + \left(\text{token} + 0.5\right) \cdot \frac{x_{\max} - x_{\min}}{V}
\end{equation}

\paragraph{Mu-law Encoding} applies logarithmic compression before binning, giving finer resolution near the center of the range:
\begin{equation}
    f(x) = \text{sign}(x) \cdot \frac{\ln(1 + \mu |x|)}{\ln(1 + \mu)}
\end{equation}

\subsection{Sequence Structure}

We flatten observations and actions into a single token sequence:
\[
    \underbrace{s^0_1, s^0_2, s^0_3, s^0_4, s^0_5, s^1_1, \ldots, s^1_5}_{\text{observations: } T_o \times 5 = 10},
    \underbrace{a^0_1, a^0_2, a^1_1, \ldots, a^{15}_2}_{\text{actions: } T_a \times 2 = 32}
\]

During training, all tokens are provided (teacher forcing). During inference, action tokens are replaced with a learned mask embedding.

\subsection{Dimension-Aware Positional Encoding}

Standard positional encodings treat the flattened sequence as a 1D list, losing the multi-dimensional structure. We propose encoding three aspects:
\begin{align}
    \text{PE}(i) &= \text{TimeEmb}(t_i) + \text{DimEmb}(d_i) + \text{TypeEmb}(\tau_i)
\end{align}
where $t_i$ is the timestep index, $d_i$ is the dimension index within that timestep, and $\tau_i \in \{0, 1\}$ indicates observation vs.\ action.

\section{Experimental Setup}

\begin{itemize}[noitemsep]
    \item \textbf{Hardware}: NVIDIA A100-PCIE-40GB
    \item \textbf{Training}: 500 epochs, batch size 256, AdamW ($\beta_1=0.9, \beta_2=0.999$)
    \item \textbf{Learning rate}: $10^{-4}$ with linear warmup (10 epochs) and cosine decay
    \item \textbf{Evaluation}: 50 episodes every 50 epochs
    \item \textbf{Logging}: Weights \& Biases
\end{itemize}

\section{Results}

\subsection{Baseline Comparison}

\begin{table}[h]
\centering
\caption{Baseline results on PushT. Diffusion Policy score from original paper.}
\begin{tabular}{lccccc}
\toprule
\textbf{Model} & \textbf{Params} & \textbf{Val Loss} & \textbf{Val Acc} & \textbf{Mean Score} & \textbf{Max Score} \\
\midrule
Diffusion Policy (DDPM) & 2.5M & --- & --- & 0.75 & --- \\
TRM V1 (tokenized, h=256) & 1.4M & 0.001 & 100.0\% & 0.061 & 0.067 \\
HRM V1 (tokenized, h=256) & 2.5M & 0.001 & 100.0\% & 0.045 & 0.059 \\
\midrule
TRM V2 (cont.\ obs, h=256) & 1.5M & 3.12 & 16.5\% & 0.107 & 0.154 \\
HRM V2 (cont.\ obs, h=256) & 2.8M & 3.22 & 15.5\% & 0.093 & 0.152 \\
\midrule
TRM V3 (regression, h=256) & 1.5M & 0.047 & --- & 0.170 & 0.270 \\
\textbf{HRM V3 (regression, h=256)} & \textbf{2.8M} & \textbf{0.035} & --- & \textbf{0.236} & \textbf{0.344} \\
TRM V3 large (h=512) & 7.5M & 0.066 & --- & 0.140 & 0.207 \\
TRM V3 deep (H5L6, ah=2) & 2.2M & 0.111 & --- & 0.091 & 0.143 \\
\bottomrule
\end{tabular}
\end{table}

\subsection{Key Findings}

\paragraph{Discrete tokenization fails for control.} V1 models achieve 100\% token prediction accuracy on training data but only 4--6\% mean coverage in evaluation. This is because: (1) small quantization errors compound over the 16-step action horizon; (2) the model memorizes exact training sequences rather than learning a generalizable policy; (3) any deviation from training observations leads to out-of-distribution states.

\paragraph{HRM outperforms TRM.} The hierarchical separation of H-level (planning) and L-level (execution) modules in HRM consistently outperforms the weight-shared TRM. At epoch 150, HRM achieves 23.6\% mean score vs. TRM's 17.0\%. This suggests that separate modules can specialize: H-level for trajectory planning and L-level for precise action computation.

\paragraph{Regression output is essential.} Moving from classification (cross-entropy over bins) to regression (MSE on continuous actions) improved scores by $3{-}4\times$. Continuous outputs avoid quantization errors and allow smoother action trajectories.

\paragraph{Observation noise augmentation prevents overfitting.} Adding Gaussian noise ($\sigma=2.0$) to observations during training significantly improved generalization to unseen initial states.

\section{Analysis}

\subsection{Recursion as Iterative Refinement}

The recursive reasoning mechanism in TRM/HRM is analogous to the denoising process in diffusion models:
\begin{itemize}[noitemsep]
    \item \textbf{Diffusion}: iteratively denoise a random Gaussian sample toward an action sequence
    \item \textbf{TRM/HRM}: iteratively refine latent states $z_H, z_L$ to produce better action predictions
\end{itemize}

The key advantage is weight sharing across all refinement steps, yielding compact models (2.8M params) compared to diffusion models (2.5M+ for state-based Diffusion Policy).

\subsection{Gap to Diffusion Policy}

Our best model (HRM V3, 0.34 max score) still significantly underperforms Diffusion Policy (0.75 mean score). Potential reasons:
\begin{itemize}[noitemsep]
    \item Diffusion Policy uses iterative denoising with hundreds of steps; we use 3--5 recursion cycles
    \item Diffusion Policy's action chunk prediction naturally handles multi-modality
    \item Our evaluation uses randomized initial states which may differ substantially from training distribution
    \item Training data is small (206 demonstrations) --- more data augmentation may help
\end{itemize}

\section{Conclusion}

We demonstrated that recursive reasoning models (TRM/HRM) can be adapted for continuous robotic control, with the key insight that continuous observation encoding and regression outputs are essential. The hierarchical HRM architecture outperforms the flat TRM, achieving 34.4\% maximum coverage on PushT. While still below Diffusion Policy's 75\%, this represents a promising direction for compact, parameter-efficient robotic policies based on iterative refinement.

Future work: (1) more recursion cycles with gradient checkpointing, (2) multi-modal output heads (GMM), (3) larger training datasets, (4) temporal action ensemble from multiple recursion depths.

\bibliographystyle{plain}
\begin{thebibliography}{9}
\bibitem{chi2023diffusion} Chi, C., et al. ``Diffusion Policy: Visuomotor Policy Learning via Action Diffusion.'' RSS 2023.
\bibitem{trm2024} Samsung SAIL Montreal. ``Tiny Recursive Models.'' 2024.
\bibitem{hrm2024} Sapient Inc. ``Hierarchical Reasoning Model.'' 2024.
\end{thebibliography}

\end{document}
